\documentclass{article}
\usepackage[portuguese]{babel}
\usepackage{listings}
\usepackage{xcolor}

\definecolor{codegreen}{rgb}{0,0.6,0}
\definecolor{codegray}{rgb}{0.5,0.5,0.5}
\definecolor{codepurple}{rgb}{0.8,0,0.2}
\definecolor{backcolour}{rgb}{.8,.8,1}

\lstdefinestyle{mystyle}{
    backgroundcolor=\color{backcolour},   
    commentstyle=\color{codegreen},
    keywordstyle=\color{blue},
    numberstyle=\tiny\color{codegray},
    stringstyle=\color{codepurple},
    basicstyle=\ttfamily\footnotesize,
    breakatwhitespace=false,         
    breaklines=true,                 
    captionpos=b,                    
    keepspaces=true,                 
    numbers=left,                    
    numbersep=5pt,                  
    showspaces=false,                
    showstringspaces=false,
    showtabs=false,                  
    tabsize=2
}

\lstset{style=mystyle}


\begin{document}

\section{Soma dos algorismos de um numero}

Comecamos salvando o numero em uma variavel \textit{n}. Criamos tambem a variavel \textit{soma} e uma variavel para iteracao \textit{i} que vai iniciar com o mesmo valor de \textit{n}.

    \begin{lstlisting}[language=Python]
    n = int(input("Digite um numero"))
    soma = 0
    i = n\end{lstlisting}

Agora vamos usar um laco while para iterar sobre cada digito de \textit{n}.

\begin{lstlisting}[language=Python]
     while(i != 0):
        aux = int(i/10)
        digito = i - (aux*10)
        soma += digito
        i = aux\end{lstlisting}

A cada iteracao, dividimos o valor de \textit{i} por 10 ignorando a parte decimal(usando \textit{int()}) e salvamos esse valor na variavel auxiliar \textit{aux}, obtemos o ultimo digito de nosso numero da subtracao $i - 10\textit{aux}$, dai entao somamos o valor do digito ao valor de \textit{soma}, e depois salvamos o valor de \textit{aux} em \textit{i} para a proxima iteracap.

\end{document}

